\begin{theorem}[11.2]\label{11.2}
    $H(\C)$, $A$~--- самосопряжённый оператор.
    \[
    \overline{\im A_{\lambda}} \oplus_{\perp} \ker A_{\lambda} = H
    \]
\end{theorem}

\begin{proof}
    По теореме \nameref{9.2} $\overline{\im A_\lambda} \oplus \ker (A_\lambda)^*=H$.
    Вынесем сопряжение: $(A - \lambda I)^* = A^* - \overline{\lambda} I$.
    % проблема: вылезет предательская лямбда с чертой
    Рассмотрим два варианта:
    \begin{enumerate}
        \item[$\lambda \in \R\!: $]
        $\lambda = \overline{\lambda}$, тогда формула автоматически верна.
        \item[$\lambda \not\in \R\!: $]
        $\overline{\lambda} \neq \lambda$. Отсюда по теореме \nameref{11.1} (пункт 2) $\lambda$ и $\overline{\lambda}$
        не являются собственными значениями, а значит, ядро тривиально.
        \end{enumerate}
\end{proof}